\section{Introduction}
\label{sec:intro}

Anyone who has used a language-model agent for a long research task knows the moment. The agent returns a tidy answer, the answer is incomplete, and the user types ``keep going''. Then ``are you sure that is the same man?'', then ``what about the ship he sailed on?'', then ``did you ever get a reply from the state library?''. The author researched 65 of his own relatives this way with Claude. The agent found real records, but only because a person kept supervising it. By the end, the definition of a thorough search lived in the supervisor's head and in his messages, not in the system. Coming back later with relative 66, there was no way to know whether the agent would check what it had checked for relative 12. A redacted excerpt and counts of those messages by type are not in this version, because the transcripts have not been exported and coded (H5, \cref{sec:results}).

Genealogy makes the problem sharper than most research, for two reasons. First, the records are dark data in the sense of \citet{heidorn2008dark} and \citet{hand2020dark}: they exist, but the systems that should find them cannot see them. A regimental history names a soldier on page 14 of a scanned book with no name index. A town clerk holds a marriage register and answers by mail. A church register, a ship's crew list, or a box of letters in a cousin's closet may be the only source. Second, the work is long. A letter to an archive can take three weeks to answer, and following up on it is part of the research. A chat window cannot wait three weeks, and a script that crashes on day four loses its place.

This paper turns the supervisor's loop into a system and proposes a way to measure it. Our claim is falsifiable: \emph{at equal tool-call budget, a durable \CR agent finds more verifiable hidden records of deceased people than single-pass, ReAct and checklist-only agents, without lowering claim precision.} We built the system and the benchmark, preregistered the comparison (\cref{sec:bench}), and ran it. The claim is not supported: at 100 tool calls a ReAct agent found more than the preregistered CR on both evaluations, and CR's advantages were precision and durability, not recall (\cref{sec:corpus,sec:results}).

\paragraph{Contributions.}
\begin{enumerate}
\item \textbf{\CR (CR)}, a procedure that turns the human ``keep digging'' loop into an explicit question graph. Each new question must be grounded in a specific prior finding, is deduplicated against a persisted question memory, is scored for novelty and expected yield, and is expanded under breadth, depth, node and time budgets with explicit stopping rules (\cref{sec:loop}).
\item \textbf{\bench}, a benchmark for long-horizon archival research on deceased people, with hidden ground-truth records stratified by darkness tier, an adjudication protocol, cost accounting, and fault injection: worker kills, flaky archives, multi-day waits and delayed archive replies (\cref{sec:bench}).
\item \textbf{A durable, auditable reference system}: CR plus heritage circuits, run on Temporal (retries, durable timers, signals for human approval and redirection, full event history), with evidence discipline (claim tiers, every published sentence tied to claim ids). It is released as open code, and the paper reports where durability and memory matter and where they do not (\cref{sec:solution,sec:results}).
\end{enumerate}

\noindent We are careful about what we do not claim. A name found in OCR text is a lead, not an identification. In the pilot, most record-level facts come from curated evaluation fixtures, labelled as stubs; live sources add leads. The generator in the current system is rule-based. All of this is stated where it applies and collected in \cref{sec:limits}.
